\chapter{v0.0.6 Frozen Study}
\label{chap:results}

This chapter preserves the frozen 149-repository v0.0.6 study and its
v0.0.5 paired comparator. The newer 207-repository v0.0.7 run, semantic
validation, and numerical benchmark are evaluated in
Chapter~\ref{chap:v007}; counts from the two populations are not pooled.

\section{Population and Execution Health}
All 149 repositories completed in v0.0.6. The frozen manifest contains 64,876
YAML files and 4,254,017 lines; 27,709 files and 1,178,908 lines satisfied the
Ansible selection criteria. AnsibleGuard and GitLeaks completed on all
repositories. ansible-lint completed on 128 and was partial on 21. The 200
diagnostics comprise 189 ansible-lint diagnostics and 11 YAML parse
diagnostics; they are execution evidence rather than code smells.

\begin{table}[H]
\centering
\caption{Frozen v0.0.6 population and execution outcomes.}
\begin{tabular}{lr}
\toprule
Measure & Value \\
\midrule
Repositories completed & 149 / 149 \\
Manifest YAML files / LOC & 64,876 / 4,254,017 \\
Selected Ansible files / LOC & 27,709 / 1,178,908 \\
Mean selected LOC per file & 42.55 \\
ansible-lint completed / partial & 128 / 21 \\
Diagnostics & 200 \\
\bottomrule
\end{tabular}
\end{table}

\section{Raw Findings and Canonical Evidence}
v0.0.6 produced 328,385 raw provider findings and 307,951 canonical
occurrences. The difference of 20,434 (6.22\% of raw findings) reflects
canonical evidence merging and does not mean that 20,434 findings were false.
The canonical density is 261.22 occurrences per KLOC, 11.11 per selected file,
and a mean of 2,066.79 per repository. These are descriptive workload measures,
not prevalence estimates for the wider Ansible ecosystem.

\begin{table}[H]
\centering
\caption{Provider findings and canonical leadership.}
\begin{tabular}{lrr}
\toprule
Provider & Raw findings & Canonical lead occurrences \\
\midrule
AnsibleGuard & 212,461 & 208,508 \\
ansible-lint & 114,866 & 98,394 \\
GitLeaks & 1,058 & 1,049 \\
\bottomrule
\end{tabular}
\end{table}

\begin{figure}[H]
\centering
\begin{tikzpicture}
\begin{axis}[ybar,width=13.5cm,height=7cm,ylabel={Findings / occurrences},
 symbolic x coords={AnsibleGuard,ansible-lint,GitLeaks},xtick=data,
 ymajorgrids=true,bar width=17pt,legend style={at={(0.5,-0.17)},anchor=north,
 legend columns=2},scaled y ticks=false,
 yticklabel style={/pgf/number format/fixed,/pgf/number format/1000 sep={,}}]
\addplot[fill=ansblue!75] coordinates {(AnsibleGuard,212461)
 (ansible-lint,114866) (GitLeaks,1058)};
\addplot[fill=ansorange!75] coordinates {(AnsibleGuard,208508)
 (ansible-lint,98394) (GitLeaks,1049)};
\legend{Raw provider findings,Canonical lead occurrences}
\end{axis}
\end{tikzpicture}
\caption{Raw contribution and canonical leadership by provider. GitLeaks is
small on this linear scale; its exact values remain visible in the table.}
\label{fig:provider_results}
\end{figure}

\section{Semantic Status and Severity}
\begin{figure}[H]
\centering
\begin{tikzpicture}
\begin{axis}[ybar,width=14cm,height=7cm,ylabel={Canonical occurrences},
 symbolic x coords={CONFIRMED,UNCERTAIN,EXPOSED,UNRESOLVED,ORDERED},
 xtick=data,x tick label style={rotate=25,anchor=east},nodes near coords,
 nodes near coords style={font=\scriptsize},bar width=16pt,ymin=0]
\addplot[fill=ansblue!70] coordinates {
 (CONFIRMED,257850) (UNCERTAIN,43669) (EXPOSED,4626)
 (UNRESOLVED,1684) (ORDERED,122)};
\end{axis}
\end{tikzpicture}
\caption{Semantic-state distribution. Small specialised states are reported
numerically because a pie chart would visually erase them.}
\end{figure}

\code{CONFIRMED} is the default state for a satisfied detector predicate, not
a manual true-positive label. \code{UNCERTAIN} records positive evidence plus
a named limitation. \code{EXPOSED} is used for credential-log analysis where
sensitive data is used and \code{no_log} is not known true.
\code{UNRESOLVED} denotes a variable use without a visible definition, guard,
producer, or uncertainty source. \code{ORDERED_CONFLICT} denotes a use before a
later modelled runtime definition.

Severity totals are 184,323 warnings (59.85\%), 116,729 notes (37.91\%), and
6,899 errors (2.24\%). Severity controls policy impact; diagnostics independently
control reactions to parser and provider health.

\begin{figure}[H]
\centering
\begin{tikzpicture}
\begin{axis}[xbar,width=12.5cm,height=5.7cm,xlabel={Canonical occurrences},
 symbolic y coords={Error,Note,Warning},ytick=data,xmajorgrids=true,
 nodes near coords,nodes near coords style={font=\scriptsize},
 xmin=0,xtick={0,50000,100000,150000,200000},scaled x ticks=false,
 xticklabel style={/pgf/number format/fixed,/pgf/number format/1000 sep={,}}]
\addplot[fill=ansblue!70] coordinates {(6899,Error) (116729,Note)
 (184323,Warning)};
\end{axis}
\end{tikzpicture}
\caption{Severity distribution. Error is intentionally retained on the same
linear scale to show its small share rather than exaggerating it.}
\label{fig:severity_results}
\end{figure}

\section{Category and Concept Distribution}
\begin{table}[H]
\centering
\caption{Canonical occurrences by category.}
\small
\begin{tabular}{lrr@{\hspace{8mm}}lrr}
\toprule
Category & Count & \% & Category & Count & \% \\
\midrule
Style & 127,420 & 41.38 & Security & 10,938 & 3.55 \\
Correctness & 56,653 & 18.40 & Documentation & 3,928 & 1.28 \\
Lint & 35,869 & 11.65 & Reliability & 3,667 & 1.19 \\
Compatibility & 34,104 & 11.07 & Testing & 1,771 & 0.58 \\
Maintainability & 31,727 & 10.30 & Structure & 1,417 & 0.46 \\
Dependencies & 457 & 0.15 & & & \\
\bottomrule
\end{tabular}
\end{table}

\begin{figure}[H]
\centering
\begin{tikzpicture}
\begin{axis}[xbar,width=13.5cm,height=9cm,xlabel={Canonical occurrences},
 symbolic y coords={Dependencies,Structure,Testing,Reliability,Documentation,
 Security,Maintainability,Compatibility,Lint,Correctness,Style},ytick=data,
 xmajorgrids=true,bar width=7pt,scaled x ticks=false,
 xticklabel style={/pgf/number format/fixed,/pgf/number format/1000 sep={,}}]
\addplot[fill=ansblue!70] coordinates {(457,Dependencies) (1417,Structure)
 (1771,Testing) (3667,Reliability) (3928,Documentation) (10938,Security)
 (31727,Maintainability) (34104,Compatibility) (35869,Lint)
 (56653,Correctness) (127420,Style)};
\end{axis}
\end{tikzpicture}
\caption{Category composition of the canonical result set. The chart measures
rule composition in this corpus, not category defect prevalence.}
\label{fig:category_results}
\end{figure}

The five dominant concepts are \code{missing-tags} (83,446),
\code{unguarded-variable} (43,035), ansible-lint issues (35,869), FQCN
(27,913), and YAML formatting (16,124). Together they account for 67.02\% of
canonical occurrences. This concentration means that any headline total must
be interpreted through rule composition, not as 307,951 independent severe
defects.

\section{Controlled Comparator: v0.0.5--v0.0.6}
v0.0.6 is the primary evaluated artifact. v0.0.5 appears only as a controlled
comparator needed to attribute semantic-engine changes. The comparison uses
the same 149 repositories, 27,709 selected files, and 1,178,908 selected LOC;
it is not a second prevalence study.
\begin{table}[H]
\centering
\caption{Paired frozen-population comparison.}
\begin{tabular}{lrrr}
\toprule
Version & Raw & Canonical & Uncertain / non-uncertain \\
\midrule
v0.0.5 & 331,123 & 297,557 & 19,431 / 278,126 \\
v0.0.6 & 328,385 & 307,951 & 43,669 / 264,282 \\
\bottomrule
\end{tabular}
\end{table}

\begin{figure}[H]
\centering
\begin{tikzpicture}
\begin{axis}[ybar,width=13cm,height=7cm,ylabel={Occurrences},
 symbolic x coords={Raw,Canonical,Uncertain,Non-uncertain},xtick=data,
 x tick label style={rotate=15,anchor=east},bar width=14pt,ymajorgrids=true,
 legend style={at={(0.5,-0.2)},anchor=north,legend columns=2},
 scaled y ticks=false,
 yticklabel style={/pgf/number format/fixed,/pgf/number format/1000 sep={,}}]
\addplot[fill=ansgray!65] coordinates {(Raw,331123) (Canonical,297557)
 (Uncertain,19431) (Non-uncertain,278126)};
\addplot[fill=ansblue!75] coordinates {(Raw,328385) (Canonical,307951)
 (Uncertain,43669) (Non-uncertain,264282)};
\legend{v0.0.5,v0.0.6}
\end{axis}
\end{tikzpicture}
\caption{Paired version behavior on identical repositories, files and LOC.
The four bars are different measures and must not be summed.}
\label{fig:paired_versions}
\end{figure}

Canonical occurrences rise while raw findings fall because v0.0.6 restores
safer occurrence identities, retains repeated assignment evidence, and moves
cases with identifiable dynamic/runtime limitations into \code{UNCERTAIN}.
Therefore the canonical delta cannot be interpreted as a simple regression or
precision change.

\section{Unguarded-Variable Mechanism Analysis}
Of 43,035 canonical \code{unguarded-variable} occurrences, 41,351 are
\code{UNCERTAIN} and 1,684 are \code{UNRESOLVED}. Five repositories dominate
the strict unresolved set:
\begin{table}[H]
\centering
\caption{Largest concentrations of unresolved variable uses.}
\begin{tabular}{lrp{7cm}}
\toprule
Repository group & Count & Dominant mechanism \\
\midrule
SAP & 630 & Dynamic operating-system/version \code{include_vars} \\
Linuxfabrik & 323 & Shared role pre-task and role-entry variables \\
Infrared & 165 & Composed variables and plugin inputs \\
Forklift & 134 & Documented command-line extra variables \\
CMS & 62 & Derived and external variables \\
\bottomrule
\end{tabular}
\end{table}

\begin{figure}[H]
\centering
\begin{tikzpicture}
\begin{axis}[xbar,width=12.5cm,height=6.5cm,xlabel={Unresolved occurrences},
 symbolic y coords={CMS,Forklift,Infrared,Linuxfabrik,SAP},ytick=data,
 xmajorgrids=true,nodes near coords,nodes near coords style={font=\scriptsize},
 xmin=0]
\addplot[fill=ansorange!75] coordinates {(62,CMS) (134,Forklift)
 (165,Infrared) (323,Linuxfabrik) (630,SAP)};
\end{axis}
\end{tikzpicture}
\caption{Concentration of strict unresolved-variable occurrences in the five
dominant repository groups.}
\label{fig:unresolved_concentration}
\end{figure}
The first four account for 74.3\% and all five for 78.0\% of unresolved
occurrences. This is evidence of a small number of missing provenance
mechanisms, not evidence that all 1,684 are runtime failures.

\section{Manual Anomaly Samples}
All 1,753 \code{set_fact} findings are \code{UNCERTAIN}. The RHEL8-CIS change
from 13 to 32 reflects restored identity for repeated assignments; sampled
flags were assignment-only. Dell samples contained template or cross-file
consumers, supporting uncertainty rather than unconditional removal.

Eleven shell findings disappeared because the source explicitly used an
ansible-lint \code{noqa} suppression for the
\code{command-instead-of-module} rule. Source sampling confirmed that the
suppression, rather than a parser loss, caused the change.

Everpure contributes 295 canonical GitLeaks occurrences. Of these, 268 repeat
one UUID-shaped example token across 61 files, primarily module examples and
documentation. All 1,049 canonical GitLeaks occurrences have
\code{secret_context=uncertain_secret}; semantic \code{CONFIRMED} means the
provider pattern matched, not that the value is an active credential.

\section{Threats to Validity}
\textbf{Construct validity:} static predicates and semantic states are not
human defect labels. \textbf{Internal validity:} paired SHAs/files/LOC reduce
population drift, but provider versions and environments must still be recorded.
\textbf{External validity:} 149 frozen repositories do not represent every
private or industrial Ansible codebase. \textbf{Conclusion validity:} no
precision or recall inference is made without a published labelled benchmark.
